\documentclass[a4paper,12pt]{article}
\usepackage[utf8]{inputenc}
\usepackage[T1]{fontenc}
\usepackage[italian]{babel}
\usepackage{amsmath,amssymb}
\usepackage{siunitx}
\usepackage{geometry}
\usepackage{graphicx}
\geometry{margin=2.5cm}

\begin{document}
\section*{Dinamica rotazionale - principio di conservazione del momento angolare}
   \begin{figure}[h]
       \centering
       \includegraphics[width=1\linewidth]{../../immagini/Esercizio59Pag310FTE.jpeg}
       \caption{Esercizio 59 pag. 310}
       \label{fig:placeholder}
   \end{figure}
 

\section*{Dati}

\[
m_p = \SI{7.50}{g}
\qquad
v_p = \SI{390}{m/s}
\qquad
l_s = \SI{0.40}{m}
\qquad
l_p = \SI{0.30}{m}
\qquad
\omega = \SI{3.20}{rad/s}
\]
\section*{Obiettivo}
La massa $M$ dello sportello
\section*{Analisi della situazione}

Consideriamo uno sportello che, in seguito all'urto con un proiettile,
inizia a ruotare in senso antiorario. Poiché il proiettile resta conficcato
nello sportello, l'urto è completamente anelastico.

Durante l'urto il sistema può essere considerato isolato rispetto
all'asse di rotazione, perché il momento delle forze esterne è trascurabile.
Per questo motivo possiamo applicare la conservazione del momento angolare
prima e dopo l'urto.

\section*{Momento angolare prima dell'urto}

Prima dell'urto, il momento angolare del proiettile rispetto all'asse di rotazione è

\[
L_1 = m_p \, v_p \, l_p
\]

\section*{Momento angolare dopo l'urto}

Dopo l'urto, il sistema formato da sportello e proiettile ruota con velocità
angolare $\omega$, quindi il momento angolare finale è

\[
L_2 = I_{\text{tot}} \, \omega
\]

dove il momento di inerzia totale vale

\[
I_{\text{tot}} = I_{\text{sportello}} + I_{\text{proiettile}}
\]

Assumendo lo sportello come un'asta che ruota attorno a un estremo, si ha

\[
I_{\text{sportello}} = \frac{1}{3} M l_s^2
\]

mentre per il proiettile, considerato come punto materiale posto a distanza
$l_p$ dall'asse di rotazione,

\[
I_{\text{proiettile}} = m_p l_p^2
\]

Quindi:

\[
L_2 = \left( \frac{1}{3} M l_s^2 + m_p l_p^2 \right)\omega
\]

\section*{Conservazione del momento angolare}

Per la conservazione del momento angolare, si ha

\[
m_p \, v_p \, l_p
=
\left( \frac{1}{3} M l_s^2 + m_p l_p^2 \right)\omega
\]

Sviluppando i calcoli:

\[
\frac{1}{3} M l_s^2 \omega
=
m_p v_p l_p - m_p l_p^2 \omega
\]

da cui si ricava

\[
M
=
\frac{3 m_p l_p \left( v_p - l_p \omega \right)}{l_s^2 \omega}
\]

\section*{Sostituzione dei valori numerici}

\[
M
=
\frac{
3 \cdot 7.50 \cdot 10^{-3} \cdot 0.30 \cdot
\left(390 - 0.30 \cdot 3.20\right)
}{
(0.40)^2 \cdot 3.20
}
\]

\[
M \approx \SI{5.12}{kg}
\]

\section*{Risultato}

La massa dello sportello è quindi

\[
\boxed{M \approx \SI{5.12}{kg}}
\]

\end{document}